% ============================================================
% Chapter 7: Conclusion — 本文（Phase 2, 2026-09-07）
% 骨子版（Phase 1, P1-P2）は git 履歴を参照。
% 構成: 何をしたか → 何が分かったか（貢献 3 点に対応）→ 限界と今後の課題 3 方向。
% 数値は第 6 章と同じ一次データ（% source: を参照）。第 1 段の再設計は A1 で実施済みのため
% 将来課題から外し、第 2 段落の介入として記した（ギャップ分析 W3）。
% 設計書: docs/superpowers/specs/2026-09-07-w3-conclusion-abstract-design.md
% 2026-10-12 改訂（8/31 報告への加藤先生の注釈の反映）: 結果を最良の重みの baseline との差で書き直し、
%   「閉じた系」「閉包率がオラクルより高い」を撤回（停止条件の充足率が高いのは規則の定義による）。
%   LLM 評価と等価な式の扱いを今後の課題に加えた（日本語版報告 §11 と同じ）。
% ============================================================

This thesis formulated one step of automated physical model building as a retrieval problem and
addressed it with a method that evaluates candidate equations as a set. Given a description of the
desired model and its input and output variables, the method presents, from a database of 11,146
equations drawn from 361 sources, the equations likely to serve in the model, and the engineer
assembles and checks the model from them. Because an equation that contains a variable other than
the inputs and outputs is useful only together with an equation that determines that variable, the
method scores each candidate not only by its relevance to the query but by its relation to the
equations already selected: a first stage narrows the database by text similarity and variable
overlap, and a second stage reranks the candidates with a multilayer perceptron that reads three
set-aware features --- complementarity, coherence, and domain agreement --- beside seven
query--equation features. A learned greedy variant selects one equation at a time, and a
self-terminating criterion, DoF-stop, decides how many equations to present: it halts once the
selected equations cover every output variable with no more unknowns than equations.

The thesis makes three contributions. First, it provides a benchmark on which a failure can be
attributed to the method: all 1,823 evaluation ground truths meet the first three requirements of
Kato and Kano, 1,000 of the 2,838 cases are synthesized by a rule that grows a ground-truth set
from an ordinary differential equation and gives each equation an output of its own, and a
stratified source-disjoint split with 10 seeds measures generalization to unseen sources. Second,
the set-aware two-stage method raises Recall@K over classical retrieval at its best weights from
0.607 to 0.743 in the mixed setting and from 0.559 to 0.689 in the synthesized-only setting, on all
10 splits, and to 0.780 and 0.740 with a variable-weighted first stage; BM25 and dense retrieval with
E5 stay at or below 0.602 and 0.555 even at their best mixing weights, and direct generation by a
large language model reaches 0.23--0.30. Third, the thesis explains the gain rather than only
reporting it. Ablation and permutation importance agree that the variable overlap between an
equation and the query --- and, among the set-aware features, the two that read shared variables
--- carries the improvement, whereas topical similarity, informative over the whole database,
becomes redundant once the first stage has screened the candidates. Two predictions drawn from this
analysis were tested by intervention and confirmed: learned greedy selection recovered the
saturation of the static reference set on the cases with the most equations ($+0.056$ at eight or
more ground-truth equations), and a variable-weighted first stage raised the candidate coverage and,
with it, the end metric. Without being told the answer size, DoF-stop stopped at the same
exact-match rate as selection that is told it, although the set at which it stops need not be
solvable.
% source: experiments/external_baselines_{A,B}.json tfidf__none__w40-60（0.6071 / 0.5585）;
%         strat_A.json / strat_B.json（0.7429 / 0.6890）; kato_0831_followup_stats.json vs_best_baseline
%         （wins 10）, stage1_B_w30-70（0.7401）, gt_checks; stage1_weight_stats.json（0.7799）;
%         external_baselines_stats.json best_per_scorer（A ≤ 0.6020, B ≤ 0.5548）;
%         llm_set_A/B/full_equiv_results.json（0.2258–0.3003）; feature_x_difficulty_stats.json /
%         pfi_results.json（変数系 +0.678 / 話題系 +0.006）; greedy_3way_stats.json
%         tests_hard_X8plus（+0.056）; dof_stop_stats.json set_exact_dof = set_exact_oracleK

Four directions follow from the limitations of Section~\ref{sec:limitations}. First, learned
greedy selection still trains with teacher forcing, which scheduled sampling addresses. The ranking
errors that remain after the two interventions lie outside every current feature: an identical
equation from another document, which an evaluation that counts equivalent equations as correct
would no longer penalize, and an equation that determines the wrong variable, which a feature
encoding the determined variable addresses. The comparison between equations, which now reads only their variable sets, could also borrow the
tree-based and embedding representations of formula search (Section~\ref{sec:rw-mathir}). DoF-stop,
which selects greedily and only counts unknowns against equations, could instead search the
candidates for combinations that meet all four requirements, as the combination method
does\cite{CombineEq}, and check that each unknown can be matched to an equation of its own. Second,
the comparison with language models should cover every case, repeat the evaluation, check the
equivalence judgments against human judgments, and include a model that reads the source documents.
Third, the evaluation is database-internal and the variable notation of the queries matches the
database; validation on real modeling tasks, including requests whose equations are absent from
the database or written in a different notation, will show how far the results carry. Fourth, the
method should be integrated into the AutoPMoB pipeline, downstream of equation extraction and
variable-equivalence judgment and upstream of a simulator, so that the engineer can assemble the
presented equations into a model, solve it, and compare it with plant data --- the test that
automated model building must finally pass.
